In this subsection, we address the decision problem formulated in Section \ref{sec:Problem formulation}.
%The best response of the sender, as given in Definition \ref{defn:Best response}, is defined in terms of the expected payoff of the sender. However, in the gameplay, the sender commits to a strategy after observing the realized feature vector $x$, which is sampled according to a probability law $D \in \fk{D}$. Consequently, instead of maximizing the expectation of his payoff, the sender could maximize his payoff function pointwise for each $x \in \cg{X}$. The next lemma states this observation.
%
%\begin{lem} Let $D \in \fk{D}$ be probability law on $\cg{X}$ and let $r$ be a receiver strategy. Then $s \in B(r)$ (best response set) if and only if
%\[s(x) := \arg\max_{s' \in \msf{S}} U(r(s'(x))) - C(s'(x),x)\; \text{almost surely}.\]
%\end{lem}
%
%\begin{proof}
%Let $\dot{\psi}$ be the probability density function of the probability law $D$. Let $s \in B(r)$. Let $Z \subset \cg{X}$ be such that for $x \in Z$, we have $s(x) \neq \arg\max_{s' \in \msf{S}} \left( U(r(s'(x))) - C(s'(x),x) \right)$. Suppose $D(Z) \neq 0$. Now, consider an alternative sender strategy $\bar{s}$ defined as
%\[\bar{s}(x) = \begin{cases}
%	s(x) & \text{for} x \in \cg{X}-Z \\
%	\arg\max_{s' \in \msf{S}} U(r(s'(x))) - C(s'(x),x) & \text{for} x \in Z.\\
%	\end{cases}
%	\]
%It then follows that
%\[
%\int_Z \mathbb{U}(x;r,s) \dot{\psi}(x) \, dx < \int_Z \mathbb{U}(x;r,\bar{s}) \dot{\psi}(x) \, dx
%\]
%and
%\[\int_{\cg{X} \setminus Z} \mathbb{U}(x;r,s) \dot{\psi}(x) \, dx = \int_{\cg{X} \setminus Z} \mathbb{U}(x;r,\bar{s}) \dot{\psi}(x) \, dx.\]
%Hence, we conclude that $\cg{U}(r,s) < \cg{U}(r,\bar{s})$. However, this contradicts the assumption that $s \in B(r)$. Therefore, we must have $D(Z) = 0$.
%Conversely, let $Z \subset \cg{X}$ such that $D(Z) = 0$, and let $\bar{s}$ be such that $\bar{s}(x) = \arg\max_{s' \in \msf{S}} \left( U(r(s'(x))) - C(s'(x),x) \right)$ for $x \in Z$. Let $s \in B(r)$. By the definition of $\bar{s}$, we have:
%\[
%\int_Z \mathbb{U}(x;r,s) \dot{\psi}(x) \, dx \leq \int_Z \mathbb{U}(x;r,\bar{s}) \dot{\psi}(x) \, dx
%\]
%and
%\[\int_{\cg{X} \setminus Z} \mathbb{U}(x;r,s) \dot{\psi}(x) \, dx = \int_{\cg{X} \setminus Z} \mathbb{U}(x;r,\bar{s}) \dot{\psi}(x) \, dx = 0.\]
%Thus, $\cg{U}(r,s) \leq \cg{U}(r,\bar{s})$. Since $s \in B(r)$, the only way to avoid a contradiction is if
%\[
%\int_Z \mathbb{U}(x;r,s) \dot{\psi}(x) \, dx = \int_Z \mathbb{U}(x;r,\bar{s}) \dot{\psi}(x) \, dx.
%\]
%This implies that the sender obtains the same payoff with $\bar{s}$ as with $s$. Therefore, $\bar{s} \in B(r)$.
%\end{proof}
We analyze the problem for various cases of utility functions.\\
%I. \textbf{$\um, \up \leq 0$:} \todo{make subsubsection instead of manual labels}\\
\subsubsection{$\um, \up \leq 0$}
We show that this case is uninteresting as the objectives of the sender and the receiver are aligned as seen in the next proposition.


\begin{thm} Let $\up,\um \leq 0$. Then,
$r(x) \equiv h(x)$ is a Stackelberg equilibrium.
	
%	 \hltodo{$h$ is a Stackelberg equilibrium.}{strategy for receiver? write $r=h$}
\end{thm}
\begin{proof}
	First, we will show that $s_h(x) = x$. Let $x \leq \mathbf{h}$, then $h(x) = \minus$.
    Since $\up \leq 0$, it follows that
\[U(h(x),x) = 0 \geq \max_{x' \in \cg{X}}U(h(x'),x) - |\phi(x')-\phi(x)|.\]
Therefore $s_h(x) = x$. Similar reasoning holds for $x > \mathbf{h}$.
	Now, the set of feature vectors whose label is correctly classified by $h$ is given by the set
	$$E(h,s_h;h) = \{x \in \cg{X}| h(s_h(x)) = h(x)\} = \cg{X}.$$
	Since $h$ correctly classifies all the feature vectors, it is a Stackelberg equilibrium.
\end{proof}


%\noindent II.  \textbf{$\up \geq 1 \geq \um$:}\\
\subsubsection{$\up \geq 1 \geq \um$}

For this case we show that any receiver $r$ whose range contains the $\plus$ label, can at most only correctly classify  feature vectors whose true labels is $\plus$.



\begin{thm}\label{prop:constant is equib}
The receiver strategy $r(c) \equiv c$ for all $x \in \cg{X}$, where $c$ is the majority label is a Stackelberg equilibrium.
	% Let $r$ be the receiver strategy such that their exists $a \in  \cg{X}$ such that
	% $r(a) = \plus$,
	% then $E(r) \supset [0,\textbf{h}]$.
\end{thm}

\begin{proof}
Let $r$ be the receiver Stackelberg equilibrium.
Then, either there does not exist any $x \in \cg{X}$ such that $r(x) = \plus$, or there exists some $t \in \cg{X}$ such that $r(t) = \plus$.
Suppose there does not exist any $x \in \cg{X}$ such that $r(x) = \plus$. Then $r(s_r(x)) \equiv \minus$ for all $x \in \cg{X}$. Thus, $E(r) \subset [\h,1]$. However, the receiver strategy $r'(x) \equiv \minus$ for all $x \in \cg{X}$ yields $E(r') = [\h,1]$. Since $r$ is a Stackelberg equilibrium, it follows that $\E(r) = [\h,1]$.
Now, suppose there exists $t \in \cg{X}$ such that $r(t) = \plus$. Since $\up \geq 1$, it follows that for all $x \in [0,\h)$, $r(s_r(x)) = \plus$. Therefore, $E(r) = [0,\h)$.
Comparing these two cases, it follows that $r(x) \equiv c$ for all $x \in \cg{X}$, where $c$ is the majority label in the Stackelberg equilibrium.
\end{proof}


%\noindent III.  \textbf{$\um,\up \geq 1$:\\}
\subsubsection{$\um,\up \geq 1$}
For this case, we show that for any receiver strategy $r$ whose range contains both $\minus$ and $\plus$ labels does not correctly classify any feature vector.

\begin{thm}
The receiver strategy $r(x) \equiv c$ for all $x \in \cg{X}$, where $c$ is the majority label is a Stackelberg equilibrium.
	
\end{thm}

\begin{proof}
Let $r$ be a receiver Stackelberg equilibrium such that there exist $a, b \in \cg{X}$ with $r(a) = \plus$ and $r(b) = \minus$. Since $\up > 1$, for all $x \in [0, \h)$, the worst-case best response is $s_r(x) = a$, so $r(s_r(x)) = \plus$. Similarly, for $x \in [\mathbf{h}, 1)$, $s_r(x) = b$, so $r(s_r(x)) = \minus$. Therefore, $E(r) = \cg{X}$ and $\cg{E}(r) = 1$.
However, for a receiver strategy $r'(x) \equiv c$ for all $x \in \cg{X}$, where $c$ is the majority label, we have $\cg{E}(r') < 0.5$. This contradicts the optimality of the Stackelberg equilibrium strategy $r$. Therefore, no such $a, b$ exist. This implies that $r$ must be a constant label strategy. Thus, $r(x) \equiv c$ for all $x \in \cg{X}$, where $c$ is the majority label, is a receiver Stackelberg equilibrium strategy.
\end{proof}


\subsubsection{$1 \geq \up \geq 0$ and  \textbf{$|\um| < \up$}}
This case includes scenarios where the sender has a strictly non-uniform preference($\um \geq 0$) as well as scenarios where the sender has a uniform preference ($-\up \leq \um < 0$).
We show that in each of these cases, there exists a receiver Stackelberg equilibrium strategy that belongs to the family of strategies $\sr{R}$ defined as follows.

Let $\ib,\jb$ be such that $0 \leq \ib \leq \h < \jb \leq 1$ and $\phi(\jb) - \phi(\ib) = \up$. Since $\phi$ is a continuously increasing function, such $\ib, \jb$ exist. Let $r(\cdot;\ib,\jb)$ denote a receiver strategy defined as follows
	\begin{equation*}\label{eqn:simple strategies}
		r(x;\ib,\jb) := \begin{cases}
			h(x) & \text{if} x \in [0,\ib] \cup [\jb,1] \\
			\Delta & \text{if} x \in (\ib,\jb)
		\end{cases}.
	\end{equation*}
Let $$\sr{R}_1 := \{r(\cdot;\ib,\jb): 0 \leq \ib \leq \h < \jb \leq 1, \phi(\jb) - \phi(\ib) = \up\}.$$
The strategy $r(\cdot,\ib,\jb) \in \sr{R}_1$ is identical to the true function $h$ except in the interval $(\ib,\jb)$, where it maps the feature vectors to label $\Delta$.
Let$$\sr{R}_2 := \{r(\cdot): r(\cdot) \equiv \minus\} \cup \{r(\cdot): r(\cdot) \equiv \plus\},$$  be the set of constant-label receiver strategies. Let $\sr{R} := \sr{R}_1 \cup \sr{R}_2$.


We define dominance between receiver strategies as follows: for a given probability law $D$ on $\cg{X}$, a strategy $r$ dominates another strategy $\rt$ if $\cg{E}(r) \leq \cg{E}(\rt)$. The next theorem shows that the family of receiver strategies $\sr{R}$ dominates every other receiver strategy. Hence, a receiver can restrict its search for a Stackelberg equilibrium to the set $\sr{R}$ itself.
\begin{thm}\label{thm:Deterministic Stackelberg equilibrium strategy} Let $D \in \fk{D}$ be a law on $\cg{X}$ and let $h(x)$ be a true function. Let $\phi$ be a cost kernel and $\up > 0$. Then their exists a  Stackelberg equilibrium strategy $r(\cdot) \in \sr{R}$.
%	\begin{enumerate}[label=(\alph*)]
%		\item If $\up < \min\{\phi(\h),1-\phi(\h)\}$, then there exist $\ib,\jb \in (0,1)$ such that $r(x;\ib,\jb) \in \sr{R}_1$ is a Stackelberg equilibrium.
%		\item If $\up \geq \min\{\phi(\h),1-\phi(\h)\}$, then $r \in \sr{R}_2$, $r(\cdot) \equiv c$, where $c$ is the majority label, is a Stackelberg equilibrium.
%	\end{enumerate}
\end{thm}



\begin{proof} The proof is given in the Proof \ref{proof:Best response r1} in the Appendix.
\end{proof}



% \begin{proof} The proof is given in the Proof \ref{proof:Deterministic domination of a family of strategies} in the Appendix.
% \end{proof}



%\noindent \textbf{V.} \textbf{$1 \geq \up \geq 0$} and \textbf{$\um \leq -\up$:}\\
\subsubsection{$1 \geq \up \geq 0$} and \textbf{$\um \leq -\up$}

This case corresponds to the sender having an uniform preference. The next theorem specifies a Stackelberg equilibrium.

\begin{thm}\label{thm:u_ < - u_+}
	Let $D \in \fk{D}$ be a probability law on $\cg{X}$ and let $h(x)$ be a true function. Let $\jb \in \cg{X}$ be such that $\phi(\jb) - \phi(\h) = \up$. Then the receiver strategy
	\[ r(x;\h,\jb):=
	\begin{cases}
		\minus & \text{for } x \in [0,\h],\\
		\Delta & \text{for } x \in (\h,\jb),\\
		\plus & \text{for } x \in [\jb,1],
	\end{cases}\]
	is a Stackelberg equilibrium.
\end{thm}

\begin{proof} The proof is given in the Proof \ref{proof:u_ < - u_+} in the Appendix.
	\end{proof}